\documentclass[12pt]{fiche-rd-info}
\usepackage[utf8]{inputenc}
\usepackage[T1]{fontenc}

\begin{document}

\authorA{Matis}{Hadet}
\authorB{Floran}{Martel}

\begin{fichesuivi}{19 octobre 2026}{24 octobre 2026}
	\tempstravailA{9}{00}
	\tempstravailB{9}{00}

	\begin{travaileffectue}
		\begin{itemize}
			\item tâche 1 : Lecture des documents fournis par Claudia Marinica.
			\item tâche 2 : Recherche de papiers sur l'apprentissage contrastif.
			\item tâche 2 : Échange avec Claudia pour exprimer notre compréhension du sujet.
			\item tâche 3 : Planification des tâches de la semaine prochaine.
			\item tâche 4 : 
		\end{itemize}
	\end{travaileffectue}

	\begin{travailnoneffectue}
		\begin{itemize}
			\item Aucune tâche non effectuée.
		\end{itemize}
	\end{travailnoneffectue}

	\begin{echange}
		\begin{itemize}
			\item L'échange avec Claudia Marinica a permis de mettre une methode de travailles en place.
		\end{itemize}
	\end{echange}

	\begin{planification}
		\begin{itemize}
			\item Trouver d'autres applications de l'apprentissage contrastif et d'autres papiers.
			\item Reflechir sur : Pourquoi ne pas utiliser certaines approches
			\item Faire un tableau de comparaison des approches -> avec des critères pour les approches et pour G-STIM
			\item Analyser l'etat de l'art de STKGPLM
			\item Regarder papier recents pour etat de l'art
		\end{itemize}
	\end{planification}
\end{fichesuivi}

\end{document}
